% SymbolMap.tex — reference sheet for the ScottDomains formalization.
% Two tables:
%   A. how pdftotext garbles the Type-3 fonts of Gunter & Scott, "Semantic
%      Domains" (1990)  ->  correct Unicode  (for READING the paper);
%   B. domain-theory symbol  <->  Unicode  <->  Lean 4 / VS Code input  <->
%      Lean display          (for READING/WRITING the Lean).
% Compile: scripts/tex2pdf.sh <this> ScottDomains/docs   (XeLaTeX + unicode-math)
\documentclass[10pt]{article}

\usepackage[letterpaper,margin=1.6cm]{geometry}
\usepackage{fontspec}
\usepackage{unicode-math}
\usepackage{booktabs}
\usepackage{longtable}
\usepackage{array}
\usepackage[dvipsnames]{xcolor}
\usepackage{newunicodechar}
% STIX Two Text lacks these arrows in text mode; route literals through math.
\newunicodechar{→}{\ensuremath{\to}}
\newunicodechar{↔}{\ensuremath{\leftrightarrow}}

% Fonts that actually cover the domain-theory glyphs (no tofu).
\setmainfont{STIX Two Text}   % text; pairs with STIX Two Math
\setmonofont{Menlo}[Scale=0.9]
\setmathfont{STIX Two Math}   % covers ⊑ ⊒ ⨆ ⨅ ≪ ↝ 𝟙 ⊗ ⊕ …

\newcommand{\li}[1]{\texttt{\textbackslash #1}}      % Lean input: \seq
\newcommand{\ex}[1]{\texttt{#1}}                      % raw extraction snippet
\newcommand{\uni}[1]{$#1$}                            % glyph via unicode-math

\setlength{\parindent}{0pt}
\renewcommand{\arraystretch}{1.25}

\begin{document}

\begin{center}
  {\Large\bfseries ScottDomains — Symbol Map}\\[2pt]
  {\footnotesize Gunter \& Scott, ``Semantic Domains'' (HTCS Vol.~B, 1990),
   pp.~633–674 — reading the paper and reading/writing the Lean.}
\end{center}

\vspace{4pt}

%==============================================================================
\section*{Table A — PDF garbling → Unicode (for reading the paper)}
%==============================================================================
{\footnotesize
The paper's PDF uses 1990-era Type-3 bitmap fonts; \texttt{pdftotext} maps many
glyphs to unrelated ASCII and silently drops the \texttt{fi\,/\,fl\,/\,ff\,/\,ffi}
ligatures. Several distinct operators all collapse to a blank gap — context (sets
vs.\ domains) disambiguates. Snippets below are quoted verbatim from the
extraction.

\begin{longtable}{@{}p{2.2cm}p{1.7cm}p{12.3cm}@{}}
\toprule
\textbf{Extraction} & \textbf{Correct} & \textbf{Meaning / example (extraction → intended)} \\
\midrule
\endfirsthead
\toprule
\textbf{Extraction} & \textbf{Correct} & \textbf{Meaning / example (extraction → intended)} \\
\midrule
\endhead
\ex{!}        & \uni{\to}      & continuous-map / function-space arrow: \ex{f : D ! D} → $f : D \to D$ \\
\ex{!}        & \uni{\omega}   & the ordinal $\omega$ (SAME glyph as the arrow): \ex{ordinal !}, \ex{cpo ! >} → $\omega^{\top}$ \\
\ex{v}        & \uni{\sqsubseteq} & the partial order (reflexive relation): \ex{binary relation v}, \ex{x v y} \\
\ex{w}        & \uni{\sqsupseteq} & the reverse order: \ex{x w y} → $x \sqsupseteq y$ \\
\ex{F}        & \uni{\bigsqcup} & directed least upper bound (join): \ex{x v F M} → $x \sqsubseteq \bigsqcup M$ \\
\ex{2}        & \uni{\in}      & set membership: \ex{x 2 D} → $x \in D$ \\
\ex{8}        & \uni{\forall}  & universal quantifier: \ex{8x 2 u} → $\forall x \in u$ \\
\ex{9}        & \uni{\exists}  & existential quantifier: \ex{9y 2 E} → $\exists y \in E$ \\
\ex{>}        & \uni{\top}     & top element: \ex{> and ?} → $\top$ and $\bot$ \\
\ex{?}        & \uni{\bot}     & bottom / least element: \ex{?D} → $\bot_{D}$, \ex{? v >} → $\bot \sqsubseteq \top$ \\
(blank gap)   & \uni{\subseteq} & subset (of sets): \ex{M  D is directed} → $M \subseteq D$; \ex{u  N} \\
(blank gap)   & \uni{\times}   & Cartesian product (of domains): \ex{D  E} → $D \times E$; \ex{fst : D  E ! D} \\
(blank gap)   & \uni{\otimes}  & smash / tensor product: \ex{smash : D  E ! D E} → $\mathrm{smash}:D\times E\to D\otimes E$; \ex{f g} → $f\otimes g$ \\
(blank gap)   & \uni{\circ}    & function composition: \ex{g = f  unsmash} → $g = f\circ\mathrm{unsmash}$; \ex{smash  (f  g)} \\
(blank gap)   & \uni{\cdot}    & multiplication / application dot: \ex{x  x} → $x\cdot x$ \\
\ex{7!}       & \uni{\mapsto}  & maps-to: \ex{x 7! x  x} → $x \mapsto x\cdot x$ \\
(vanished)    & \uni{\lambda}  & Church's binder drops entirely: \ex{x: x  x} → $\lambda x.\,x\cdot x$; \ex{-notation} → $\lambda$-notation \\
line-lead \ex{=} & \uni{\cong} & isomorphism ($\cong$; the tilde is dropped, leaving a bare \ex{=} at line start): \ex{D <newline> = D ! D} → $D \cong (D \to D)$ \\
\ex{\{}       & \uni{\text{--}} & en-dash in number/page ranges: \ex{633\{674} → 633–674 \\
\ex{h \ldots i} & \uni{\langle\,\ldots\,\rangle} & tuple / pairing angle brackets: \ex{hf1, : : :, fni} → $\langle f_1,\ldots,f_n\rangle$ \\
\ex{: : :}    & \uni{\ldots}   & ellipsis (also the dotted leaders in the table of contents): \ex{(D1; : : :; Dn)} → $(D_1,\ldots,D_n)$ \\
\ex{fi} drop  & fi & \ex{de nition}→definition, \ex{ nite}→finite, \ex{ xed}→fixed, \ex{bi nite}→bifinite, \ex{ nitary}→finitary, \ex{satis es}→satisfies, \ex{signi cant}→significant, \ex{ rst}→first, \ex{ nd}→find \\
\ex{fl} drop  & fl & \ex{re exive}→reflexive, \ex{brie y}→briefly \\
\ex{ff} drop  & ff & \ex{di erent}→different, \ex{e ects}→effects, \ex{e ectively}→effectively \\
\ex{ffi} drop & ffi & \ex{diculty}→difficulty, \ex{dicult}→difficult \\
\bottomrule
\end{longtable}
}

\vspace{4pt}

%==============================================================================
\section*{Table B — Symbol ↔ Unicode ↔ Lean input ↔ Lean display (for the Lean)}
%==============================================================================
{\footnotesize
Lean input = the \texttt{\textbackslash}-sequence typed in the VS Code
\texttt{lean4} extension (equivalently \texttt{lean4-mode} in Emacs), from the
extension's \texttt{abbreviations.json}. The primary canonical abbreviation is
listed first; common aliases follow. ``Lean shows'' is the glyph the editor
substitutes — identical to the Symbol column, rendered here through
\texttt{unicode-math}. Every sequence below was checked against the installed
abbreviation table; none are unverified.

\begin{longtable}{@{}p{1.1cm}p{6.9cm}p{5.6cm}p{1.4cm}@{}}
\toprule
\textbf{Symbol} & \textbf{Unicode (code point + name)} & \textbf{Lean input} & \textbf{Lean shows} \\
\midrule
\endfirsthead
\toprule
\textbf{Symbol} & \textbf{Unicode (code point + name)} & \textbf{Lean input} & \textbf{Lean shows} \\
\midrule
\endhead
\uni{\sqsubseteq} & U+2291 SQUARE IMAGE OF OR EQUAL TO      & \li{sqsubseteq} \; (\li{squb})              & \uni{\sqsubseteq} \\
\uni{\sqsupseteq} & U+2292 SQUARE ORIGINAL OF OR EQUAL TO   & \li{sqsupseteq} \; (\li{squp})              & \uni{\sqsupseteq} \\
\uni{\sqcup}      & U+2294 SQUARE CUP                       & \li{sqcup} \; (\li{lub}, \li{join})         & \uni{\sqcup} \\
\uni{\sqcap}      & U+2293 SQUARE CAP                       & \li{sqcap} \; (\li{glb}, \li{meet})         & \uni{\sqcap} \\
\uni{\bigsqcup}   & U+2A06 N-ARY SQUARE UNION OPERATOR      & \li{bigsqcup} \; (\li{Sqcup}, \li{Sup}, \li{supr}) & \uni{\bigsqcup} \\
\uni{\bigsqcap}   & U+2A05 N-ARY SQUARE INTERSECTION OP.    & \li{bigsqcap} \; (\li{Sqcap}, \li{Inf}, \li{infi}) & \uni{\bigsqcap} \\
\uni{\bot}        & U+22A5 UP TACK                          & \li{bot}                                    & \uni{\bot} \\
\uni{\top}        & U+22A4 DOWN TACK                        & \li{top}                                    & \uni{\top} \\
\uni{\ll}         & U+226A MUCH LESS-THAN (way-below)       & \li{ll}                                     & \uni{\ll} \\
\uni{\to}         & U+2192 RIGHTWARDS ARROW                 & \li{to} \; (\li{r}, \texttt{\textbackslash->}, \li{rightarrow}) & \uni{\to} \\
\uni{\mapsto}     & U+21A6 RIGHTWARDS ARROW FROM BAR        & \li{mapsto} \; (\texttt{\textbackslash|->}) & \uni{\mapsto} \\
\uni{\circ}       & U+2218 RING OPERATOR (composition)      & \li{circ} \; (\li{o}, \li{comp})            & \uni{\circ} \\
\uni{\times}      & U+00D7 MULTIPLICATION SIGN              & \li{times} \; (\li{x})                      & \uni{\times} \\
\uni{\otimes}     & U+2297 CIRCLED TIMES (smash / tensor)   & \li{otimes} \; (\li{ox})                    & \uni{\otimes} \\
\uni{\oplus}      & U+2295 CIRCLED PLUS (sum)               & \li{oplus} \; (\texttt{\textbackslash o+})  & \uni{\oplus} \\
\uni{\cong}       & U+2245 APPROXIMATELY EQUAL TO (iso)     & \li{cong} \; (\li{iso}, \texttt{\textbackslash\textasciitilde=}) & \uni{\cong} \\
\uni{\simeq}      & U+2243 ASYMPTOTICALLY EQUAL TO (equiv)  & \li{simeq} \; (\li{equiv})                  & \uni{\simeq} \\
\uni{\infty}      & U+221E INFINITY                         & \li{infty}                                  & \uni{\infty} \\
\uni{\omega}      & U+03C9 GREEK SMALL LETTER OMEGA         & \li{omega} \; (\li{om})                     & \uni{\omega} \\
\uni{\beta}       & U+03B2 GREEK SMALL LETTER BETA          & \li{beta} \; (\li{b})                       & \uni{\beta} \\
\uni{\lambda}     & U+03BB GREEK SMALL LETTER LAMDA         & \li{lambda} \; (\li{la}, \li{fun})          & \uni{\lambda} \\
\uni{\alpha}      & U+03B1 GREEK SMALL LETTER ALPHA         & \li{alpha} \; (\li{a})                      & \uni{\alpha} \\
\uni{\Gamma}      & U+0393 GREEK CAPITAL LETTER GAMMA       & \li{Gamma} \; (\li{G})                      & \uni{\Gamma} \\
\uni{\Pi}         & U+03A0 GREEK CAPITAL LETTER PI          & \li{Pi} \; (\li{P})                         & \uni{\Pi} \\
\uni{\Sigma}      & U+03A3 GREEK CAPITAL LETTER SIGMA       & \li{Sigma} \; (\li{S})                      & \uni{\Sigma} \\
\uni{\in}         & U+2208 ELEMENT OF                       & \li{in} \; (\li{mem})                       & \uni{\in} \\
\uni{\notin}      & U+2209 NOT AN ELEMENT OF                & \li{notin} \; (\li{nin})                    & \uni{\notin} \\
\uni{\forall}     & U+2200 FOR ALL                         & \li{forall} \; (\li{all})                   & \uni{\forall} \\
\uni{\exists}     & U+2203 THERE EXISTS                     & \li{exists} \; (\li{ex})                    & \uni{\exists} \\
\uni{\subseteq}   & U+2286 SUBSET OF OR EQUAL TO            & \li{subseteq} \; (\li{ss}, \li{sub})        & \uni{\subseteq} \\
\uni{\subset}     & U+2282 SUBSET OF                        & \li{ssubset} \; (\li{ssub}) \; [\li{subset}$\Rightarrow\subseteq$] & \uni{\subset} \\
\uni{\wedge}      & U+2227 LOGICAL AND                      & \li{and} \; (\li{wedge})                    & \uni{\wedge} \\
\uni{\vee}        & U+2228 LOGICAL OR                       & \li{or} \; (\li{vee})                       & \uni{\vee} \\
\uni{\lnot}       & U+00AC NOT SIGN                         & \li{not} \; (\li{neg}, \li{lnot})           & \uni{\lnot} \\
\uni{\le}         & U+2264 LESS-THAN OR EQUAL TO            & \li{le} \; (\li{leq})                       & \uni{\le} \\
\uni{\ge}         & U+2265 GREATER-THAN OR EQUAL TO         & \li{ge} \; (\li{geq})                       & \uni{\ge} \\
\uni{↝} & U+219D RIGHTWARDS WAVE ARROW       & \li{leadsto} \; (\texttt{\textbackslash r\textasciitilde}) & \uni{↝} \\
\uni{𝟙}   & U+1D7D9 MATH DOUBLE-STRUCK DIGIT ONE    & \li{b1}                                     & \uni{𝟙} \\
\bottomrule
\end{longtable}
}

\vspace{2pt}
{\scriptsize Sources: extraction via \texttt{pdftotext} on
\texttt{ScottDomains/papers/Gunter Scott 1990.pdf}; Lean input from the
\texttt{lean4} abbreviation table (\texttt{abbreviations.json}).}

\end{document}
